% Generated mechanically from the frozen spaces-topologies.tex target.
% Do not edit this generated file; edit the bound locale source instead.

% OK, start here.
%

\setcounter{chapter}{72}
\chapter{代数空间上的拓扑}


\phantomsection
\label{spaces-topologies-section-phantom}





\section{引言}
\label{spaces-topologies-section-introduction}

\noindent
本章介绍代数空间范畴上的若干拓扑。可与 \cite{SGA1}、
\cite{Ner}、\cite{LM-B} 及 \cite{Kn} 中的材料比较。
在此之前需要指出，可以选取许多不同的位点（定义见《位点》定义
\ref{sites-definition-site}），而它们在所依范畴上给出相同的层概念。
因此，我们的选择可能与参考文献中的选择略有不同，但最终得到相同的上同调群等。




\section{一般构造过程}
\label{spaces-topologies-section-procedure}

\noindent
本节说明构造下文所用位点的一般过程。读者若未读过《拓扑》第
\ref{topologies-section-procedure} 节，则这段讨论几乎无法理解。

\medskip\noindent
设 $S$ 为基概形。从 $S$ 以及任意一个希望包含的 $S$-概形集合出发，
按《集合》引理 \ref{sets-lemma-construct-category} 构造任一范畴
$\Sch_\alpha$。再从范畴 $\Sch_\alpha$ 与 fppf 覆盖类出发，
按《集合》引理 \ref{sets-lemma-coverings-site} 选取
$\Sch_\alpha$ 上任一覆盖集合 $\text{Cov}_{fppf}$。
以 $\Sch_{fppf}$ 表示所得大 fppf 位点，以
$(\Sch/S)_{fppf}$ 表示相应的 $S$ 的大 fppf 位点。
（以上完全依照《拓扑》第 \ref{topologies-section-fppf} 节的规定。）

\medskip\noindent
作出上述选择后，$S$ 上代数空间的范畴具有一个同构类集合。
这是因为任一 $S$ 上代数空间都可写成 $U/R$，其中
$j : R \to U \times_S U$ 是某个平展等价关系，且
$U, R \in \Ob((\Sch/S)_{fppf})$，见《空间》引理
\ref{spaces-lemma-space-presentation}。
因此，可在 $S$ 上代数空间的范畴中找出一个满子范畴
$\textit{Spaces}/S$；它的对象构成集合，并且每个代数空间都同构于
$\textit{Spaces}/S$ 的某个对象。固定这样一个范畴。

\medskip\noindent
在以下各节中，给定拓扑 $\tau$ 后，大位点
$(\textit{Spaces}/S)_\tau$（相应地，$S$ 上代数空间 $X$ 的大位点
$(\textit{Spaces}/X)_\tau$）的所依范畴是
$\textit{Spaces}/S$（相应地，是 $\textit{Spaces}/S$ 的子范畴
$\textit{Spaces}/X$，见《范畴》例
\ref{categories-example-category-over-X}）。
按通常方式把它变成位点：定义一类 $\tau$-覆盖，并应用《集合》引理
\ref{sets-lemma-coverings-site} 选取一个足够大的覆盖集合来定义该拓扑。

\medskip\noindent
需要指出，{\it 代数空间 $X$ 的小平展位点 $X_\etale$} 已在
《空间的性质》定义
\ref{spaces-properties-definition-etale-site} 中定义。
它的对象是在 $X$ 上平展的概形；由代数空间的定义，这样的对象很多。
% P09-ERR-0312
不过，从本章的视角看（比较《拓扑》定义
\ref{topologies-definition-big-small-etale}），更自然的是
《空间的性质》定义
\ref{spaces-properties-definition-spaces-etale-site} 中的位点
$X_{spaces, \etale}$。这两个位点定义同一拓扑斯，见《空间的性质》引理
\ref{spaces-properties-lemma-compare-etale-sites}。
本章不再重新定义它们，而直接使用。






\section{Zariski 拓扑}
\label{spaces-topologies-section-zariski}

\noindent
在《空间》第 \ref{spaces-section-Zariski} 节中，我们介绍了用开子空间
覆盖代数空间的 Zariski 覆盖概念。将开子空间替换为开浸入，得到如下相应概念。

\begin{definition}
\label{spaces-topologies-definition-zariski-covering}
设 $S$ 为概形，$X$ 为 $S$ 上的代数空间。
{\it $X$ 的 Zariski 覆盖}是 $S$ 上代数空间的一个态射族
$\{f_i : X_i \to X\}_{i \in I}$，其中每个 $f_i$ 都是开浸入，并且
$$
|X| = \bigcup\nolimits_{i \in I} |f_i|(|X_i|),
$$
即这些态射联合满。
\end{definition}

\noindent
尽管 Zariski 覆盖偶尔有用，但代数空间范畴上相应的拓扑确实太粗，
因而并不特别有用。尽管如此，它仍定义一个位点。

\begin{lemma}
\label{spaces-topologies-lemma-zariski}
设 $S$ 为概形，$X$ 为 $S$ 上的代数空间。
\begin{enumerate}
\item 若 $X' \to X$ 是同构，则 $\{X' \to X\}$ 是 $X$ 的 Zariski 覆盖。
\item 若 $\{X_i \to X\}_{i\in I}$ 是 Zariski 覆盖，并且对每个
$i$，$\{X_{ij} \to X_i\}_{j\in J_i}$ 是 Zariski 覆盖，则
$\{X_{ij} \to X\}_{i \in I, j\in J_i}$ 是 Zariski 覆盖。
\item 若 $\{X_i \to X\}_{i\in I}$ 是 Zariski 覆盖，而
$X' \to X$ 是代数空间的态射，则
$\{X' \times_X X_i \to X'\}_{i\in I}$ 是 Zariski 覆盖。
\end{enumerate}
\end{lemma}

\begin{proof}
略。
\end{proof}










\section{平展拓扑}
\label{spaces-topologies-section-etale}

\noindent
本节讨论代数空间的平展覆盖概念，并定义代数空间的大平展位点。
请与《拓扑》第 \ref{topologies-section-etale} 节比较。

\begin{definition}
\label{spaces-topologies-definition-etale-covering}
设 $S$ 为概形，$X$ 为 $S$ 上的代数空间。
{\it $X$ 的平展覆盖}是 $S$ 上代数空间的一个态射族
$\{f_i : X_i \to X\}_{i \in I}$，其中每个 $f_i$ 都平展，并且
$$
|X| = \bigcup\nolimits_{i \in I} |f_i|(|X_i|),
$$
即这些态射联合满。
\end{definition}

\noindent
这恰好就是《拓扑》定义 \ref{topologies-definition-etale-covering}。
特别地，若 $X$ 及所有 $X_i$ 都是概形，便得到通常的概形平展覆盖概念。

\begin{lemma}
\label{spaces-topologies-lemma-zariski-etale}
任一 Zariski 覆盖都是平展覆盖。
\end{lemma}

\begin{proof}
由定义以及开浸入是平展态射这一事实即可得出（由于浸入可表，
该事实由《态射》引理 \ref{morphisms-lemma-open-immersion-etale}
经《空间》引理
\ref{spaces-lemma-representable-transformations-property-implication} 得到）。
\end{proof}

\begin{lemma}
\label{spaces-topologies-lemma-etale}
设 $S$ 为概形，$X$ 为 $S$ 上的代数空间。
\begin{enumerate}
\item 若 $X' \to X$ 是同构，则 $\{X' \to X\}$ 是 $X$ 的平展覆盖。
\item 若 $\{X_i \to X\}_{i\in I}$ 是平展覆盖，并且对每个
$i$，$\{X_{ij} \to X_i\}_{j\in J_i}$ 是平展覆盖，则
$\{X_{ij} \to X\}_{i \in I, j\in J_i}$ 是平展覆盖。
\item 若 $\{X_i \to X\}_{i\in I}$ 是平展覆盖，而
$X' \to X$ 是代数空间的态射，则
$\{X' \times_X X_i \to X'\}_{i\in I}$ 是平展覆盖。
\end{enumerate}
\end{lemma}

\begin{proof}
略。
\end{proof}

\noindent
下一个引理说明位点
$(\textit{Spaces}/X)_\etale$ 与 $(\textit{Spaces}/X)_{smooth}$
具有相同的层范畴。

\begin{lemma}
\label{spaces-topologies-lemma-etale-dominates-smooth}
设 $S$ 为概形，$X$ 为 $S$ 上的代数空间。
设 $\{X_i \to X\}_{i \in I}$ 是 $X$ 的光滑覆盖。
则存在 $X$ 的平展覆盖 $\{U_j \to X\}_{j \in J}$，
它加细 $\{X_i \to X\}_{i \in I}$。
\end{lemma}

\begin{proof}
先取概形 $U$ 及满平展态射 $U \to X$。
对每个 $i$，取概形 $W_i$ 及满平展态射 $W_i \to X_i$。
于是 $\{W_i \to X\}_{i \in I}$ 是加细
$\{X_i \to X\}_{i \in I}$ 的光滑覆盖。因此
$\{W_i \times_X U \to U\}_{i \in I}$ 是概形的光滑覆盖。
由《态射续论》引理 \ref{more-morphisms-lemma-etale-dominates-smooth}，
可取加细 $\{W_i \times_X U \to U\}$ 的平展覆盖
$\{U_j \to U\}$。于是 $\{U_j \to X\}_{j \in J}$ 是加细
$\{X_i \to X\}_{i \in I}$ 的平展覆盖。
\end{proof}

\begin{definition}
\label{spaces-topologies-definition-big-etale-site}
设 $S$ 为概形。大平展位点 {\it $(\textit{Spaces}/S)_\etale$}
是按如下方式构造的任一位点：
\begin{enumerate}
\item 按《拓扑》第 \ref{topologies-section-etale} 节选取一个
大平展位点 $(\Sch/S)_\etale$。
\item 取 $S$ 上代数空间的范畴 $\textit{Spaces}/S$ 为所依范畴
（关于它为何是集合，见第 \ref{spaces-topologies-section-procedure} 节的讨论）。
\item 按《集合》引理 \ref{sets-lemma-coverings-site}，
从范畴 $\textit{Spaces}/S$ 及定义
\ref{spaces-topologies-definition-etale-covering} 的平展覆盖类出发，选取任一覆盖集合。
\end{enumerate}
\end{definition}

\noindent
定义完毕后，可通过局部化得到代数空间的平展位点。

\begin{definition}
\label{spaces-topologies-definition-big-small-etale}
设 $S$ 为概形，$(\textit{Spaces}/S)_\etale$ 如定义
\ref{spaces-topologies-definition-big-etale-site}。
设 $X$ 为 $S$ 上的代数空间，即 $(\textit{Spaces}/S)_\etale$
的对象。则 $X$ 的大平展位点
{\it $(\textit{Spaces}/X)_\etale$} 是《位点》第
\ref{sites-section-localize} 节所引入的位点
$(\textit{Spaces}/S)_\etale$ 在 $X$ 处的局部化。
\end{definition}

\noindent
回顾定义中的 $S$ 上代数空间 $X$，我们已经定义了小平展位点
$X_{spaces, \etale}$ 与 $X_\etale$，见《空间的性质》第
\ref{spaces-properties-section-etale-site} 节。
我们借助包含函子 $X_\etale \subset X_{spaces, \etale}$
（《空间的性质》引理
\ref{spaces-properties-lemma-compare-etale-sites}）
默认识别相应的拓扑斯，并称之为 $X$ 的小平展拓扑斯。
下面建立这些位点所关联的拓扑斯之间的一些关系。

\begin{lemma}
\label{spaces-topologies-lemma-put-in-T-etale}
设 $S$ 为概形，$f : Y \to X$ 为
$(\textit{Spaces}/S)_\etale$ 中的态射。包含函子
$Y_{spaces, \etale} \to (\textit{Spaces}/X)_\etale$
是余连续的，并诱导拓扑斯态射
$$
i_f :
\Sh(Y_\etale)
\longrightarrow
\Sh((\textit{Spaces}/X)_\etale)
$$
对 $(\textit{Spaces}/X)_\etale$ 上的层 $\mathcal{G}$，有公式
$(i_f^{-1}\mathcal{G})(U/Y) = \mathcal{G}(U/X)$。
函子 $i_f^{-1}$ 还具有左伴随 $i_{f, !}$，后者与纤维积和等化子可交换。
\end{lemma}

\begin{proof}
以 $u : Y_{spaces, \etale} \to (\textit{Spaces}/X)_\etale$ 表示该函子。
换言之，给定对应于 $Y_{spaces, \etale}$ 中一个对象的平展态射
$j : U \to Y$，置
% P09-ERR-0313
$u(U \to T) = (f \circ j : U \to S)$.
范畴 $Y_{spaces, \etale}$ 有纤维积和等化子，且 $u$ 与它们可交换。
% P09-ERR-0314
显然 $u$ 余连续。函子 $u$ 也连续，因为 $u$ 将覆盖变为覆盖并与
纤维积可交换。因此引理由《位点》引理
\ref{sites-lemma-when-shriek} 和
\ref{sites-lemma-preserve-equalizers} 得出。
\end{proof}

\begin{lemma}
\label{spaces-topologies-lemma-at-the-bottom-etale}
设 $S$ 为概形，$X$ 为 $(\textit{Spaces}/S)_\etale$ 的对象。
包含函子 $X_{spaces, \etale} \to (\textit{Spaces}/X)_\etale$
满足《位点》引理 \ref{sites-lemma-bigger-site} 的假设，因而诱导位点态射
$$
\pi_X :
(\textit{Spaces}/X)_\etale
\longrightarrow
X_{spaces, \etale}
$$
以及拓扑斯态射
$$
i_X :
\Sh(X_\etale)
\longrightarrow
\Sh((\textit{Spaces}/X)_\etale)
$$
使得 $\pi_X \circ i_X = \text{id}$。此外，
$i_X = i_{\text{id}_X}$，其中 $i_{\text{id}_X}$ 如引理
\ref{spaces-topologies-lemma-put-in-T-etale} 所述。特别地，函子
$i_X^{-1} = \pi_{X, *}$ 由规则
$i_X^{-1}(\mathcal{G})(U/X) = \mathcal{G}(U/X)$.
描述。
\end{lemma}

\begin{proof}
在此情形，函子
$u : X_{spaces, \etale} \to (\textit{Spaces}/X)_\etale$
除具有上面引理 \ref{spaces-topologies-lemma-put-in-T-etale} 的证明中所见性质外，
还是全忠实的，并将终对象变为终对象。
本引理由《位点》引理 \ref{sites-lemma-bigger-site} 得出。
\end{proof}

\begin{definition}
\label{spaces-topologies-definition-restriction-small-etale}
在引理 \ref{spaces-topologies-lemma-at-the-bottom-etale} 的情形中，函子
$i_X^{-1} = \pi_{X, *}$ 常称为{\it 限制到小平展位点}。
% P09-ERR-0315
对大平展位点上的层 $\mathcal{F}$，常以
$\mathcal{F}|_{X_\etale}$ 表示这一限制。
\end{definition}

\noindent
采用这一记号，对大位点上的层 $\mathcal{F}$ 及小位点上的层
$\mathcal{G}$，有
\begin{align*}
\Mor_{\Sh(X_\etale)}(
\mathcal{F}|_{X_\etale},
\mathcal{G})
& =
\Mor_{\Sh((\textit{Spaces}/X)_\etale)}(
\mathcal{F},
i_{X, *}\mathcal{G}) \\
\Mor_{\Sh(X_\etale)}(
\mathcal{G},
\mathcal{F}|_{X_\etale})
& =
\Mor_{\Sh((\textit{Spaces}/X)_\etale)}(
\pi_X^{-1}\mathcal{G},
\mathcal{F})
\end{align*}
此外，有 $(i_{X, *}\mathcal{G})|_{X_\etale} = \mathcal{G}$，
也有 $(\pi_X^{-1}\mathcal{G})|_{X_\etale} = \mathcal{G}$。

\begin{lemma}
\label{spaces-topologies-lemma-morphism-big-etale}
设 $S$ 为概形，$f : Y \to X$ 为
$(\textit{Spaces}/S)_\etale$ 中的态射。函子
$$
u :
(\textit{Spaces}/Y)_\etale
\longrightarrow
(\textit{Spaces}/X)_\etale,
\quad
V/Y \longmapsto V/X
$$
是余连续的，并有连续右伴随
$$
v :
(\textit{Spaces}/X)_\etale
\longrightarrow
(\textit{Spaces}/Y)_\etale,
\quad
(U \to X) \longmapsto (U \times_X Y \to Y).
$$
它们诱导同一拓扑斯态射
$$
f_{big} :
\Sh((\textit{Spaces}/Y)_\etale)
\longrightarrow
\Sh((\textit{Spaces}/X)_\etale)
$$
且有 $f_{big}^{-1}(\mathcal{G})(U/Y) = \mathcal{G}(U/X)$，
$f_{big, *}(\mathcal{F})(U/X) = \mathcal{F}(U \times_X Y/Y)$。
此外，$f_{big}^{-1}$ 有左伴随 $f_{big!}$，后者与纤维积和等化子可交换。
\end{lemma}

\begin{proof}
函子 $u$ 余连续、连续，并与纤维积及等化子可交换
（细节略去；比较引理 \ref{spaces-topologies-lemma-put-in-T-etale} 的证明）。
因此可应用《位点》引理 \ref{sites-lemma-when-shriek} 和
\ref{sites-lemma-preserve-equalizers}，得到 $f_{big}^{-1}$
的公式以及 $f_{big!}$ 的存在性。此外，函子 $v$ 是右伴随，
因为给定 $U/Y$ 与 $V/X$，有
$\Mor_X(u(U), V) = \Mor_Y(U, V \times_X Y)$，正合所需。
于是可应用《位点》引理
\ref{sites-lemma-have-functor-other-way} 和
\ref{sites-lemma-have-functor-other-way-morphism}，
得到 $f_{big, *}$ 的公式。
\end{proof}

\begin{lemma}
\label{spaces-topologies-lemma-morphism-big-small-etale}
设 $S$ 为概形，$f : Y \to X$ 为
$(\textit{Spaces}/S)_\etale$ 中的态射。
\begin{enumerate}
\item % P09-ERR-0316
有 $i_f = f_{big} \circ i_T$，其中 $i_f$ 如引理
\ref{spaces-topologies-lemma-put-in-T-etale}，$i_T$ 如引理
\ref{spaces-topologies-lemma-at-the-bottom-etale}。
\item % P09-ERR-0317
函子 $X_{spaces, \etale} \to T_{spaces, \etale}$，
$(U \to X) \mapsto (U \times_X Y \to Y)$ 连续，并诱导位点态射
$$
f_{spaces, \etale} : Y_{spaces, \etale} \longrightarrow X_{spaces, \etale}
$$
相应的小平展拓扑斯态射记作
$$
f_{small} : \Sh(Y_\etale) \to \Sh(X_\etale)
$$
有 $f_{small, *}(\mathcal{F})(U/X) = \mathcal{F}(U \times_X Y/Y)$。
\item 有如下位点态射的交换图
$$
\xymatrix{
Y_{spaces, \etale} \ar[d]_{f_{spaces, \etale}} &
(\textit{Spaces}/Y)_\etale \ar[d]^{f_{big}} \ar[l]^-{\pi_Y}\\
X_{spaces, \etale} &
(\textit{Spaces}/X)_\etale \ar[l]_-{\pi_X}
}
$$
因而作为拓扑斯态射，
$f_{small} \circ \pi_Y = \pi_X \circ f_{big}$。
\item 有 $f_{small} = \pi_X \circ f_{big} \circ i_Y = \pi_X \circ i_f$。
\end{enumerate}
\end{lemma}

\begin{proof}
% P09-ERR-0318
等式 $i_f = f_{big} \circ i_Y$ 由等式
$i_f^{-1} = i_T^{-1} \circ f_{big}^{-1}$ 得出；
后者由上述函子描述即知。因此得到 (1)。

\medskip\noindent
% P09-ERR-0319
在《空间的性质》引理
\ref{spaces-properties-lemma-functoriality-etale-site} 中已经说明，
函子 $u : X_{spaces, \etale} \to Y_{spaces, \etale}$，
$u(U \to X) = (U \times_X Y \to Y)$，给出一个位点态射及相应的
小平展拓扑斯态射。正像的描述是显然的。

\medskip\noindent
(3) 成立，因为 $\pi_X$ 和 $\pi_Y$ 由包含函子给出，而
$f_{spaces, \etale}$ 和 $f_{big}$ 由基变换函子
$U \mapsto U \times_X Y$ 给出。

\medskip\noindent
在 (3) 两端前复合 $i_Y$ 即得 (4)。
\end{proof}

\noindent
在该引理的情形中，采用定义
\ref{spaces-topologies-definition-restriction-small-etale} 的术语，
对 $Y$ 的大平展位点上的层 $\mathcal{F}$，有
$$
(f_{big, *}\mathcal{F})|_{X_\etale} =
f_{small, *}(\mathcal{F}|_{Y_\etale}),
$$
% P09-ERR-0320
该等式由引理中的位点图交换即知，因为限制到 $Y$（相应地 $X$）
的小平展位点由 $\pi_{Y, *}$（相应地 $\pi_{X, *}$）给出。
涉及拉回与限制的类似公式不成立。

\begin{lemma}
\label{spaces-topologies-lemma-composition-etale}
设 $S$ 为概形。给定 $(\textit{Spaces}/S)_\etale$ 中的态射
$f : X \to Y$、$g : Y \to Z$，有
$g_{big} \circ f_{big} = (g \circ f)_{big}$，且
$g_{small} \circ f_{small} = (g \circ f)_{small}$.
\end{lemma}

\begin{proof}
对大位点上的函子，此结论由引理 \ref{spaces-topologies-lemma-morphism-big-etale}
中正像与拉回的简单描述得出。对小位点上的函子，此结论由引理
\ref{spaces-topologies-lemma-morphism-big-small-etale} 中正像函子的描述得出。
\end{proof}

\begin{lemma}
\label{spaces-topologies-lemma-morphism-big-small-cartesian-diagram-etale}
设 $S$ 为概形。考虑 $(\textit{Spaces}/S)_\etale$ 中的笛卡尔图
$$
\xymatrix{
Y' \ar[r]_{g'} \ar[d]_{f'} & Y \ar[d]^f \\
X' \ar[r]^g & X
}
$$
则
$i_g^{-1} \circ f_{big, *} = f'_{small, *} \circ (i_{g'})^{-1}$
且 $g_{big}^{-1} \circ f_{big, *} = f'_{big, *} \circ (g'_{big})^{-1}$。
\end{lemma}

\begin{proof}
由于该图是笛卡儿的，对 $U'/X'$ 有
$U' \times_{X'} Y' = U' \times_X Y$。因此
$i_g^{-1} \circ f_{big, *}$ 与
$f'_{small, *} \circ (i_{g'})^{-1}$ 都把
$(\textit{Spaces}/Y)_\etale$ 上的层 $\mathcal{F}$ 送到
$X'_\etale$ 上的层
$U' \mapsto \mathcal{F}(U' \times_{X'} Y')$
（使用引理 \ref{spaces-topologies-lemma-put-in-T-etale} 和
\ref{spaces-topologies-lemma-morphism-big-etale}）。
第二个等式可用相同方法证明，也可由一般性的《位点》引理
\ref{sites-lemma-localize-morphism} 推出。
\end{proof}

\begin{remark}
\label{spaces-topologies-remark-change-topologies-ringed}
位点 $(\textit{Spaces}/X)_\etale$ 与 $X_{spaces, \etale}$
都带有结构层。对小平展位点，我们已在《空间的性质》第
\ref{spaces-properties-section-structure-sheaf} 节看到这一点。
大平展位点 $(\textit{Spaces}/X)_\etale$ 上的结构层
$\mathcal{O}$ 定义为：给对象 $U$ 指派 $U$ 的结构层的整体截面。
这是有意义的，因为 $U$ 本身毕竟是代数空间，因而具有结构层。
由于 $\mathcal{O}_U$ 是 $U$ 的平展位点上的层，如此定义的预层
$\mathcal{O}$ 满足 $U$ 的覆盖的层条件；也就是说，
$\mathcal{O}$ 是层。
可以把上面定义的态射 $i_f$、$\pi_X$、$i_X$、$f_{small}$ 及
$f_{big}$ 分别提升为带环位点、带环拓扑斯的态射。下面逐一处理。
\begin{enumerate}
\item 在引理 \ref{spaces-topologies-lemma-put-in-T-etale} 中，以 $\mathcal{O}$
表示 $(\textit{Spaces}/X)_\etale$ 上的结构层。
按构造有
$(i_f^{-1}\mathcal{O})(U/Y) = \mathcal{O}_U(U) = \mathcal{O}_Y(U)$。
% P09-ERR-0321
从而有同构 $i_f^\sharp : i_f^{-1}\mathcal{O} \to \mathcal{O}_Y$。
\item 在引理 \ref{spaces-topologies-lemma-at-the-bottom-etale} 中已经指出，
$i_X$ 是 $f = \text{id}_X$ 时 $i_f$ 的特例，故回到情形 (1)。
\item 在引理 \ref{spaces-topologies-lemma-at-the-bottom-etale} 中，态射
$\pi_X$ 满足 $(\pi_{X, *}\mathcal{O})(U) = \mathcal{O}(U) =
\mathcal{O}_X(U)$。因此可用它定义
$\pi_X^\sharp : \mathcal{O}_X \to \pi_{X, *}\mathcal{O}$。
\item 在引理 \ref{spaces-topologies-lemma-morphism-big-small-etale} 中，
《空间的性质》引理
\ref{spaces-properties-lemma-morphism-ringed-topoi}
讨论了把 $f_{small}$ 扩张为带环拓扑斯态射。
\item 在引理 \ref{spaces-topologies-lemma-morphism-big-small-etale} 中，
函子 $f_{big}^{-1}$ 只不过是通过包含函子
$(\textit{Spaces}/Y)_\etale \to (\textit{Spaces}/X)_\etale$
作限制。设 $\mathcal{O}_1$ 为 $(\textit{Spaces}/X)_\etale$
上的结构层，$\mathcal{O}_2$ 为 $(\textit{Spaces}/Y)_\etale$
上的结构层。得到典范同构
$f_{big}^\sharp : f_{big}^{-1}\mathcal{O}_1 \to \mathcal{O}_2$。
\end{enumerate}
此外，在这些定义下，复合也正确相容。略去详细陈述与证明。
\end{remark}









\section{光滑拓扑}
\label{spaces-topologies-section-smooth}

\noindent
本节讨论代数空间的光滑覆盖概念，并定义代数空间的大光滑位点。
请与《拓扑》第 \ref{topologies-section-smooth} 节比较。

\begin{definition}
\label{spaces-topologies-definition-smooth-covering}
设 $S$ 为概形，$X$ 为 $S$ 上的代数空间。
{\it $X$ 的光滑覆盖}是 $S$ 上代数空间的一个态射族
$\{f_i : X_i \to X\}_{i \in I}$，其中每个 $f_i$ 都光滑，并且
$$
|X| = \bigcup\nolimits_{i \in I} |f_i|(|X_i|),
$$
即这些态射联合满。
\end{definition}

\noindent
这恰好就是《拓扑》定义 \ref{topologies-definition-smooth-covering}。
特别地，若 $X$ 及所有 $X_i$ 都是概形，便得到通常的概形光滑覆盖概念。

\begin{lemma}
\label{spaces-topologies-lemma-zariski-etale-smooth}
任一平展覆盖都是光滑覆盖；因此更进一步，任一 Zariski 覆盖都是光滑覆盖。
\end{lemma}

\begin{proof}
这由定义、平展态射是光滑态射这一事实
（《空间的态射》引理 \ref{spaces-morphisms-lemma-etale-smooth}）
及引理 \ref{spaces-topologies-lemma-zariski-etale} 即得。
\end{proof}

\begin{lemma}
\label{spaces-topologies-lemma-smooth}
设 $S$ 为概形，$X$ 为 $S$ 上的代数空间。
\begin{enumerate}
\item 若 $X' \to X$ 是同构，则 $\{X' \to X\}$ 是 $X$ 的光滑覆盖。
\item 若 $\{X_i \to X\}_{i\in I}$ 是光滑覆盖，并且对每个
$i$，$\{X_{ij} \to X_i\}_{j\in J_i}$ 是光滑覆盖，则
$\{X_{ij} \to X\}_{i \in I, j\in J_i}$ 是光滑覆盖。
\item 若 $\{X_i \to X\}_{i\in I}$ 是光滑覆盖，而
$X' \to X$ 是代数空间的态射，则
$\{X' \times_X X_i \to X'\}_{i\in I}$ 是光滑覆盖。
\end{enumerate}
\end{lemma}

\begin{proof}
略。
\end{proof}

\noindent
待续……









\section{Syntomic 拓扑}
\label{spaces-topologies-section-syntomic}

\noindent
本节讨论代数空间的 syntomic 覆盖概念，并定义代数空间的大
syntomic 位点。请与《拓扑》第
\ref{topologies-section-syntomic} 节比较。

\begin{definition}
\label{spaces-topologies-definition-syntomic-covering}
设 $S$ 为概形，$X$ 为 $S$ 上的代数空间。
{\it $X$ 的 syntomic 覆盖}是 $S$ 上代数空间的一个态射族
$\{f_i : X_i \to X\}_{i \in I}$，其中每个 $f_i$ 都是 syntomic
态射，并且
$$
|X| = \bigcup\nolimits_{i \in I} |f_i|(|X_i|),
$$
即这些态射联合满。
\end{definition}

\noindent
这恰好就是《拓扑》定义 \ref{topologies-definition-syntomic-covering}。
特别地，若 $X$ 及所有 $X_i$ 都是概形，便得到通常的概形
syntomic 覆盖概念。

\begin{lemma}
\label{spaces-topologies-lemma-zariski-etale-smooth-syntomic}
任一光滑覆盖都是 syntomic 覆盖；因此更进一步，任一平展覆盖或
Zariski 覆盖都是 syntomic 覆盖。
\end{lemma}

\begin{proof}
这由定义、光滑态射是 syntomic 态射这一事实
（《空间的态射》引理
\ref{spaces-morphisms-lemma-smooth-syntomic}）
及引理 \ref{spaces-topologies-lemma-zariski-etale-smooth} 即得。
\end{proof}

\begin{lemma}
\label{spaces-topologies-lemma-syntomic}
设 $S$ 为概形，$X$ 为 $S$ 上的代数空间。
\begin{enumerate}
\item 若 $X' \to X$ 是同构，则 $\{X' \to X\}$ 是 $X$ 的
syntomic 覆盖。
\item 若 $\{X_i \to X\}_{i\in I}$ 是 syntomic 覆盖，并且对每个
$i$，$\{X_{ij} \to X_i\}_{j\in J_i}$ 是 syntomic 覆盖，则
$\{X_{ij} \to X\}_{i \in I, j\in J_i}$ 是 syntomic 覆盖。
\item 若 $\{X_i \to X\}_{i\in I}$ 是 syntomic 覆盖，而
$X' \to X$ 是代数空间的态射，则
$\{X' \times_X X_i \to X'\}_{i\in I}$ 是 syntomic 覆盖。
\end{enumerate}
\end{lemma}

\begin{proof}
略。
\end{proof}

\noindent
待续……










\section{Fppf 拓扑}
\label{spaces-topologies-section-fppf}

\noindent
本节讨论代数空间的 fppf 覆盖概念，并定义代数空间的大 fppf 位点。
请与《拓扑》第 \ref{topologies-section-fppf} 节比较。

\begin{definition}
\label{spaces-topologies-definition-fppf-covering}
设 $S$ 为概形，$X$ 为 $S$ 上的代数空间。
{\it $X$ 的 fppf 覆盖}是 $S$ 上代数空间的一个态射族
$\{f_i : X_i \to X\}_{i \in I}$，其中每个 $f_i$ 都平坦且局部有限表示，
并且
$$
|X| = \bigcup\nolimits_{i \in I} |f_i|(|X_i|),
$$
即这些态射联合满。
\end{definition}

\noindent
这恰好就是《拓扑》定义 \ref{topologies-definition-fppf-covering}。
特别地，若 $X$ 及所有 $X_i$ 都是概形，便得到通常的概形 fppf 覆盖概念。

\begin{lemma}
\label{spaces-topologies-lemma-zariski-etale-smooth-syntomic-fppf}
任一 syntomic 覆盖都是 fppf 覆盖；因此更进一步，任一光滑、平展或
Zariski 覆盖都是 fppf 覆盖。
\end{lemma}

\begin{proof}
这由定义、syntomic 态射平坦且局部有限表示这一事实
（《空间的态射》引理
\ref{spaces-morphisms-lemma-syntomic-locally-finite-presentation} 和
\ref{spaces-morphisms-lemma-syntomic-flat}）
及引理 \ref{spaces-topologies-lemma-zariski-etale-smooth-syntomic} 即得。
\end{proof}

\begin{lemma}
\label{spaces-topologies-lemma-fppf}
设 $S$ 为概形，$X$ 为 $S$ 上的代数空间。
\begin{enumerate}
\item 若 $X' \to X$ 是同构，则 $\{X' \to X\}$ 是 $X$ 的 fppf 覆盖。
\item 若 $\{X_i \to X\}_{i\in I}$ 是 fppf 覆盖，并且对每个
$i$，$\{X_{ij} \to X_i\}_{j\in J_i}$ 是 fppf 覆盖，则
$\{X_{ij} \to X\}_{i \in I, j\in J_i}$ 是 fppf 覆盖。
\item 若 $\{X_i \to X\}_{i\in I}$ 是 fppf 覆盖，而
$X' \to X$ 是代数空间的态射，则
$\{X' \times_X X_i \to X'\}_{i\in I}$ 是 fppf 覆盖。
\end{enumerate}
\end{lemma}

\begin{proof}
略。
\end{proof}

\begin{lemma}
\label{spaces-topologies-lemma-refine-fppf-schemes}
设 $S$ 为概形，$X$ 为 $S$ 上的代数空间。假设
$\mathcal{U} = \{f_i : X_i \to X\}_{i \in I}$ 是 $X$ 的 fppf 覆盖。
则存在 $\mathcal{U}$ 的加细
$\mathcal{V} = \{g_i : T_i \to X\}$，它是 fppf 覆盖且每个
$T_i$ 都是概形。
\end{lemma}

\begin{proof}
略。提示：对每个 $i$，取概形 $T_i$ 及满平展态射
$T_i \to X_i$，再验证 $\{T_i \to X\}$ 是 fppf 覆盖。
\end{proof}

\begin{lemma}
\label{spaces-topologies-lemma-fppf-covering-surjective}
设 $S$ 为概形。设 $\{f_i : X_i \to X\}_{i \in I}$
为 $S$ 上代数空间的 fppf 覆盖。则层态射
$$
\coprod X_i \longrightarrow X
$$
是满的。
\end{lemma}

\begin{proof}
这由《空间》引理
\ref{spaces-lemma-surjective-flat-locally-finite-presentation} 得出。
若对本引理的含义有疑问，另见《空间》注
\ref{spaces-remark-warning}。
\end{proof}

\begin{definition}
\label{spaces-topologies-definition-big-fppf-site}
设 $S$ 为概形。大 fppf 位点 {\it $(\textit{Spaces}/S)_{fppf}$}
是按如下方式构造的任一位点：
\begin{enumerate}
\item 按《拓扑》第 \ref{topologies-section-fppf} 节选取一个大
fppf 位点 $(\Sch/S)_{fppf}$。
\item 取 $S$ 上代数空间的范畴 $\textit{Spaces}/S$ 为所依范畴
（关于它为何是集合，见第 \ref{spaces-topologies-section-procedure} 节的讨论）。
\item 按《集合》引理 \ref{sets-lemma-coverings-site}，
从范畴 $\textit{Spaces}/S$ 及定义
\ref{spaces-topologies-definition-fppf-covering} 的 fppf 覆盖类出发，选取任一覆盖集合。
\end{enumerate}
\end{definition}

\noindent
定义完毕后，可通过局部化得到代数空间的 fppf 位点。

\begin{definition}
\label{spaces-topologies-definition-big-small-fppf}
设 $S$ 为概形，$(\textit{Spaces}/S)_{fppf}$ 如定义
\ref{spaces-topologies-definition-big-fppf-site}。
设 $X$ 为 $S$ 上的代数空间，即 $(\textit{Spaces}/S)_{fppf}$
的对象。则 $X$ 的大 fppf 位点
{\it $(\textit{Spaces}/X)_{fppf}$} 是《位点》第
\ref{sites-section-localize} 节所引入的位点
$(\textit{Spaces}/S)_{fppf}$ 在 $X$ 处的局部化。
\end{definition}

\noindent
下面建立这些位点所关联的拓扑斯之间的一些关系。

\begin{lemma}
\label{spaces-topologies-lemma-morphism-big-fppf}
设 $S$ 为概形，$f : Y \to X$ 为 $S$ 上代数空间的态射。函子
$$
u : (\textit{Spaces}/Y)_{fppf} \longrightarrow (\textit{Spaces}/X)_{fppf},
\quad
V/Y \longmapsto V/X
$$
是余连续的，并有连续右伴随
% P09-ERR-0322
$$
v : (\textit{Spaces}/X)_{fppf} \longrightarrow (\textit{Spaces}/Y)_{fppf},
\quad
(U \to Y) \longmapsto (U \times_X Y \to Y).
$$
它们诱导同一拓扑斯态射
$$
f_{big} :
\Sh((\textit{Spaces}/Y)_{fppf})
\longrightarrow
\Sh((\textit{Spaces}/X)_{fppf})
$$
且有 $f_{big}^{-1}(\mathcal{G})(U/Y) = \mathcal{G}(U/X)$，
$f_{big, *}(\mathcal{F})(U/X) = \mathcal{F}(U \times_X Y/Y)$。
此外，$f_{big}^{-1}$ 有左伴随 $f_{big!}$，后者与纤维积和等化子可交换。
\end{lemma}

\begin{proof}
函子 $u$ 余连续、连续，并与纤维积及等化子可交换。因此可应用
《位点》引理 \ref{sites-lemma-when-shriek} 和
\ref{sites-lemma-preserve-equalizers}，得到 $f_{big}^{-1}$
的公式以及 $f_{big!}$ 的存在性。此外，函子 $v$ 是右伴随，
因为给定
% P09-ERR-0323
$U/T$ 与 $V/X$，有
$\Mor_X(u(U), V) = \Mor_Y(U, V \times_X Y)$，正合所需。
于是可应用《位点》引理
\ref{sites-lemma-have-functor-other-way} 和
\ref{sites-lemma-have-functor-other-way-morphism}，
得到 $f_{big, *}$ 的公式。
\end{proof}

\begin{lemma}
\label{spaces-topologies-lemma-composition-fppf}
设 $S$ 为概形。给定 $S$ 上代数空间的态射
$f : X \to Y$、$g : Y \to Z$，有
$g_{big} \circ f_{big} = (g \circ f)_{big}$.
\end{lemma}

\begin{proof}
这由引理 \ref{spaces-topologies-lemma-morphism-big-fppf} 中大位点上函子的
正像与拉回的简单描述得出。
\end{proof}








\section{ph 拓扑}
\label{spaces-topologies-section-ph}

\noindent
本节定义 ph 拓扑。这是由平展覆盖和正常满态射生成的拓扑，见引理
\ref{spaces-topologies-lemma-characterize-sheaf}。

\begin{definition}
\label{spaces-topologies-definition-ph-covering}
设 $S$ 为概形，$X$ 为 $S$ 上的代数空间。
{\it $X$ 的 ph 覆盖}是 $S$ 上代数空间的一个态射族
$\{X_i \to X\}_{i \in I}$，
% P09-ERR-0324
使得 $f_i$ 局部有限型，并且对每个仿射 $U$ 及态射 $U \to X$，
存在加细族 $\{X_i \times_X U \to U\}_{i \in I}$ 的标准 ph 覆盖
$\{U_j \to U\}_{j = 1, \ldots, m}$。
\end{definition}

\noindent
% P09-ERR-0325
换言之，存在指标 $i_1, \ldots, i_m \in I$ 及态射
$h_j : U_j \to X_{i_j}$，使得
$f_{i_j} \circ h_j = h \circ g_j$。注意，若 $X$ 及所有 $X_i$ 均
可表，则由《拓扑》定义 \ref{topologies-definition-ph-covering}，
这与概形的 ph 覆盖相同。

\begin{lemma}
\label{spaces-topologies-lemma-zariski-etale-smooth-syntomic-fppf-ph}
任一 fppf 覆盖都是 ph 覆盖；因此更进一步，任一 syntomic、光滑、
平展或 Zariski 覆盖都是 ph 覆盖。
\end{lemma}

\begin{proof}
只需证明 fppf 覆盖是 ph 覆盖，其余结论便由引理
\ref{spaces-topologies-lemma-zariski-etale-smooth-syntomic-fppf} 得出。
设 $\{X_i \to X\}_{i \in I}$ 是基概形 $S$ 上代数空间的 fppf 覆盖。
设 $U$ 为仿射概形并给定态射 $U \to X$。可以将 fppf 覆盖
% P09-ERR-0326
$\{X_i \times_U U \to U\}_{i \in I}$ 加细为 fppf 覆盖
$\{T_i \to U\}_{i \in I}$，其中 $T_i$ 是概形
（引理 \ref{spaces-topologies-lemma-refine-fppf-schemes}）。
再由《态射续论》引理 \ref{more-morphisms-lemma-fppf-ph}
（以及概形的 ph 覆盖定义），可找出加细
$\{T_i \to U\}_{i \in I}$ 的标准 ph 覆盖
$\{U_j \to U\}_{j = 1, \ldots, m}$。
因此按定义，$\{X_i \to X\}_{i \in I}$ 是 ph 覆盖。
\end{proof}

\begin{lemma}
\label{spaces-topologies-lemma-surjective-proper-ph}
设 $S$ 为概形，$f : Y \to X$ 为 $S$ 上代数空间的固有满态射。
则 $\{Y \to X\}$ 是 ph 覆盖。
\end{lemma}

\begin{proof}
设 $U \to X$ 为态射，其中 $U$ 仿射。
由 Chow 引理（其弱形式见《空间的上同调》引理
\ref{spaces-cohomology-lemma-weak-chow}），存在概形的固有满态射
$V \to U$，它经 $Y \times_X U \to U$ 分解。
取 $V$ 的任一有限仿射开覆盖，便得到 $U$ 的一个标准 ph 覆盖，
它加细
% P09-ERR-0327
$\{X \times_Y U \to U\}$，正合所需。
\end{proof}

\begin{lemma}
\label{spaces-topologies-lemma-ph}
设 $S$ 为概形，$X$ 为 $S$ 上的代数空间。
\begin{enumerate}
\item 若 $X' \to X$ 是同构，则 $\{X' \to X\}$ 是 $X$ 的 ph 覆盖。
\item 若 $\{X_i \to X\}_{i\in I}$ 是 ph 覆盖，并且对每个
$i$，$\{X_{ij} \to X_i\}_{j\in J_i}$ 是 ph 覆盖，则
$\{X_{ij} \to X\}_{i \in I, j\in J_i}$ 是 ph 覆盖。
\item 若 $\{X_i \to X\}_{i\in I}$ 是 ph 覆盖，而
$X' \to X$ 是代数空间的态射，则
$\{X' \times_X X_i \to X'\}_{i\in I}$ 是 ph 覆盖。
\end{enumerate}
\end{lemma}

\begin{proof}
(1) 是显然的。考虑 (3) 中的 $g : X' \to X$ 及 ph 覆盖
$\{X_i \to X\}_{i\in I}$。由《空间的态射》引理
\ref{spaces-morphisms-lemma-base-change-finite-type}，态射
$X' \times_X X_i \to X'$ 局部有限型。
若 $h' : Z \to X'$ 是从仿射概形到 $X'$ 的态射，则置
$h = g \circ h' : Z \to X$。关于 $\{X_i \to X\}_{i\in I}$
的假设意味着，存在标准 ph 覆盖
$\{Z_j \to Z\}_{j = 1, \ldots, n}$ 以及盖住 $h$ 的态射
$Z_j \to X_{i(j)}$，其中某些 $i(j) \in I$。
由纤维积的泛性质，还得到 $h'$ 上的态射
$Z_j \to X' \times_X X_{i(j)}$。
因此 $\{X' \times_X X_i \to X'\}_{i\in I}$ 是 ph 覆盖，
这就证明了 (3)。

\medskip\noindent
设 $\{X_i \to X\}_{i\in I}$ 与
$\{X_{ij} \to X_i\}_{j\in J_i}$ 如 (2) 所述。
设 $h : Z \to X$ 为从仿射概形到 $X$ 的态射。
由假设，存在标准 ph 覆盖
$\{Z_j \to Z\}_{j = 1, \ldots, n}$ 及盖住 $h$ 的态射
$h_j : Z_j \to X_{i(j)}$，其中某些 $i(j) \in I$。
仍由假设，存在标准 ph 覆盖
$\{Z_{j, l} \to Z_j\}_{l = 1, \ldots, n(j)}$
以及盖住 $h_j$ 的态射 $Z_{j, l} \to X_{i(j)j(l)}$，
其中某些 $j(l) \in J_{i(j)}$。
由《拓扑》引理 \ref{topologies-lemma-refine-by-standard-ph}，
族 $\{Z_{j, l} \to Z\}$ 可被某个标准 ph 覆盖加细。
因此 $\{X_{ij} \to X\}_{i \in I, j\in J_i}$ 是 ph 覆盖。
\end{proof}

\begin{definition}
\label{spaces-topologies-definition-big-ph-site}
设 $S$ 为概形。大 ph 位点 {\it $(\textit{Spaces}/S)_{ph}$}
是按如下方式构造的任一位点：
\begin{enumerate}
\item 按《拓扑》第 \ref{topologies-section-ph} 节选取一个大 ph 位点
$(\Sch/S)_{ph}$。
\item 取 $S$ 上代数空间的范畴 $\textit{Spaces}/S$ 为所依范畴
（关于它为何是集合，见第 \ref{spaces-topologies-section-procedure} 节的讨论）。
\item 按《集合》引理 \ref{sets-lemma-coverings-site}，
从范畴 $\textit{Spaces}/S$ 及定义
\ref{spaces-topologies-definition-ph-covering} 的 ph 覆盖类出发，选取任一覆盖集合。
\end{enumerate}
\end{definition}

\noindent
定义完毕后，可通过局部化得到代数空间的 ph 位点。

\begin{definition}
\label{spaces-topologies-definition-big-small-ph}
设 $S$ 为概形，$(\textit{Spaces}/S)_{ph}$ 如定义
\ref{spaces-topologies-definition-big-ph-site}。
设 $X$ 为 $S$ 上的代数空间，即 $(\textit{Spaces}/S)_{ph}$ 的对象。
则 $X$ 的大 ph 位点 {\it $(\textit{Spaces}/X)_{ph}$} 是
《位点》第 \ref{sites-section-localize} 节所引入的位点
$(\textit{Spaces}/S)_{ph}$ 在 $X$ 处的局部化。
\end{definition}

\noindent
下面给出先前所承诺的 ph 层刻画。

\begin{lemma}
\label{spaces-topologies-lemma-characterize-sheaf}
设 $S$ 为概形，$X$ 为 $S$ 上的代数空间。
设 $\mathcal{F}$ 为 $(\textit{Spaces}/X)_{ph}$ 上的预层。
则 $\mathcal{F}$ 是层，当且仅当
\begin{enumerate}
\item $\mathcal{F}$ 对平展覆盖满足层条件；并且
\item 若 $f : V \to U$ 是 $(\textit{Spaces}/X)_{ph}$ 中的固有满态射，
则 $\mathcal{F}(U)$ 双射地映到两个映射
$\mathcal{F}(V) \to \mathcal{F}(V \times_U V)$ 的等化子。
\end{enumerate}
\end{lemma}

\begin{proof}
只需证明若 (1) 与 (2) 成立，则 $\mathcal{F}$ 是层。
设 $\{T_i \to T\}$ 为 ph 覆盖，即
$(\textit{Spaces}/X)_{ph}$ 中的覆盖。我们验证此覆盖的层条件。
设 $s_i \in \mathcal{F}(T_i)$ 为一些截面，它们限制到
$T_i \times_T T_{i'}$ 后相同。我们将证明存在唯一截面
% P09-ERR-0328
$s \in \mathcal{F}$，它在 $T_i$ 上限制为 $s_i$。
设 $\{U_j \to T\}$ 为平展覆盖，其中 $U_j$ 仿射。
由性质 (1)，为构造 $s$，只需构造截面
$s_j \in \mathcal{F}(U_j)$，使其在
% P09-ERR-0329
$U_j \cap U_{j'}$ 上相符。
考虑 ph 覆盖 $\{T_i \times_T U_j \to U_j\}$。
于是 $s_{ji} = s_i|_{T_i \times_T U_j}$ 是一些截面，它们在
$(T_i \times_T U_j) \times_{U_j} (T_{i'} \times_T U_j)$ 上相符。
取固有满态射 $V_j \to U_j$ 及有限仿射开覆盖
$V_j = \bigcup V_{jk}$，使标准 ph 覆盖
$\{V_{jk} \to U_j\}$ 加细
$\{T_i \times_T U_j \to U_j\}$。
若 $s_{jk} \in \mathcal{F}(V_{jk})$ 表示经隐含态射将
$s_{ji}$ 拉回到 $V_{jk}$ 所得的截面，则 $s_{jk}$ 黏合为截面
$s'_j \in \mathcal{F}(V_j)$。再次利用重叠处的一致性可知，
$s'_j$ 属于两个映射
$\mathcal{F}(V_j) \to \mathcal{F}(V_j \times_{U_j} V_j)$ 的等化子。
因此由 (2)，$s'_j$ 来自唯一截面
$s_j \in \mathcal{F}(U_j)$。略去验证这些截面
$s_j$ 具有全部所需性质。
\end{proof}

\noindent
下面建立这些位点所关联的拓扑斯之间的一些关系。

\begin{lemma}
\label{spaces-topologies-lemma-morphism-big-ph}
设 $S$ 为概形，$f : Y \to X$ 为 $S$ 上代数空间的态射。函子
$$
u : (\textit{Spaces}/Y)_{ph} \longrightarrow (\textit{Spaces}/X)_{ph},
\quad
V/Y \longmapsto V/X
$$
是余连续的，并有连续右伴随
% P09-ERR-0330
$$
v : (\textit{Spaces}/X)_{ph} \longrightarrow (\textit{Spaces}/Y)_{ph},
\quad
(U \to Y) \longmapsto (U \times_X Y \to Y).
$$
它们诱导同一拓扑斯态射
$$
f_{big} :
\Sh((\textit{Spaces}/Y)_{ph})
\longrightarrow
\Sh((\textit{Spaces}/X)_{ph})
$$
且有 $f_{big}^{-1}(\mathcal{G})(U/Y) = \mathcal{G}(U/X)$，
$f_{big, *}(\mathcal{F})(U/X) = \mathcal{F}(U \times_X Y/Y)$。
此外，$f_{big}^{-1}$ 有左伴随 $f_{big!}$，后者与纤维积和等化子可交换。
\end{lemma}

\begin{proof}
函子 $u$ 余连续、连续，并与纤维积及等化子可交换。因此可应用
《位点》引理 \ref{sites-lemma-when-shriek} 和
\ref{sites-lemma-preserve-equalizers}，得到 $f_{big}^{-1}$
的公式以及 $f_{big!}$ 的存在性。此外，函子 $v$ 是右伴随，
因为给定
% P09-ERR-0331
$U/T$ 与 $V/X$，有
$\Mor_X(u(U), V) = \Mor_Y(U, V \times_X Y)$，正合所需。
于是可应用《位点》引理
\ref{sites-lemma-have-functor-other-way} 和
\ref{sites-lemma-have-functor-other-way-morphism}，
得到 $f_{big, *}$ 的公式。
\end{proof}

\begin{lemma}
\label{spaces-topologies-lemma-composition-ph}
设 $S$ 为概形。给定 $S$ 上代数空间的态射
$f : X \to Y$、$g : Y \to Z$，有
$g_{big} \circ f_{big} = (g \circ f)_{big}$.
\end{lemma}

\begin{proof}
这由引理 \ref{spaces-topologies-lemma-morphism-big-ph} 中大位点上函子的
正像与拉回的简单描述得出。
\end{proof}

\begin{lemma}
\label{spaces-topologies-lemma-cech-enough}
设 $S$ 为概形，$X$ 为 $S$ 上的代数空间。
设 $P$ 为 $(\textit{Spaces}/X)_{fppf}$ 中对象的一种性质，
使得每当 $\{U_i \to U\}$ 是
$(\textit{Spaces}/X)_{fppf}$ 中的覆盖时，都有
$$
P(U_{i_0} \times_U \ldots \times_U U_{i_p})
\text{ 对所有 }
p \geq 0,\ i_0, \ldots, i_p \in I
\Rightarrow P(U)
$$
若对所有在 $X$ 上平坦、局部有限表示的仿射 $U$ 都有 $P(U)$，
则 $P(X)$。
\end{lemma}

\begin{proof}
设 $U$ 为 $X$ 上局部有限表示的分离代数空间。可选择平展覆盖
$\{U_i \to U\}_{i \in I}$，其中
% P09-ERR-0332
$V_i$ 仿射。由于 $U$ 分离，
$U_{i_0} \times_U \ldots \times_U U_{i_p}$ 总是仿射，
故总有 $P(U_{i_0} \times_U \ldots \times_U U_{i_p})$。
于是 $P(U)$ 成立。选择一个由仿射概形不交并而成的概形 $U$
以及满平展态射 $U \to X$。
则 $U \times_X \ldots \times_X U$（共有 $p + 1$ 个因子）
是 $X$ 上平展的分离代数空间。因此由上可知
$P(U \times_X \ldots \times_X U)$，从而 $P(X)$ 成立。
\end{proof}












\section{Fpqc 拓扑}
\label{spaces-topologies-section-fpqc}

\noindent
我们简要讨论代数空间的 fpqc 覆盖概念。请与《拓扑》第
\ref{topologies-section-fpqc} 节比较。
我们将在《空间的下降》命题
\ref{spaces-descent-proposition-fpqc-descent-quasi-coherent} 中证明，
% P09-ERR-0333
拟凝聚层可沿这些覆盖下降。

\begin{definition}
\label{spaces-topologies-definition-fpqc-covering}
设 $S$ 为概形，$X$ 为 $S$ 上的代数空间。
{\it $X$ 的 fpqc 覆盖}是代数空间的一个态射族
$\{f_i : X_i \to X\}_{i \in I}$，其中每个 $f_i$ 都平坦；
并且对每个仿射概形 $Z$ 及态射 $h : Z \to X$，都存在加细族
$\{X_i \times_X Z \to Z\}_{i \in I}$ 的标准 fpqc 覆盖
$\{g_j : Z_j \to Z\}_{j = 1, \ldots, m}$。
\end{definition}

\noindent
% P09-ERR-0334
换言之，存在指标 $i_1, \ldots, i_m \in I$ 及态射
$h_j : Z_j \to X_{i_j}$，使得
$f_{i_j} \circ h_j = h \circ g_j$。注意，若 $X$ 及所有 $X_i$ 均
可表，则由《拓扑》引理
\ref{topologies-lemma-fpqc-covering-affines-mapping-in}，
这与概形的 fpqc 覆盖相同。

\begin{lemma}
\label{spaces-topologies-lemma-zariski-etale-smooth-syntomic-fppf-fpqc}
任一 fppf 覆盖都是 fpqc 覆盖；因此更进一步，任一 syntomic、光滑、
平展或 Zariski 覆盖都是 fpqc 覆盖。
\end{lemma}

\begin{proof}
只需证明 fppf 覆盖是 fpqc 覆盖，其余结论便由引理
\ref{spaces-topologies-lemma-zariski-etale-smooth-syntomic-fppf} 得出。
设 $\{f_i : U_i \to U\}_{i \in I}$ 是 $S$ 上代数空间的 fppf 覆盖。
按定义，$f_i$ 平坦，这验证了定义
\ref{spaces-topologies-definition-fpqc-covering} 的第一个条件。
为验证第二个条件，设 $V \to U$ 为态射，其中 $V$ 仿射。
可选择平展覆盖 $\{V_{ij} \to V \times_U U_i\}$，其中 $V_{ij}$ 仿射。
于是复合
$f_{ij} : V_{ij} \to V \times_U U_i \to V$ 平坦且局部有限表示，
因为它们是具有这些性质的态射的复合
（《空间的态射》引理
\ref{spaces-morphisms-lemma-composition-finite-presentation},
\ref{spaces-morphisms-lemma-composition-flat},
\ref{spaces-morphisms-lemma-etale-flat} 及
\ref{spaces-morphisms-lemma-etale-locally-finite-presentation}).
因此这些态射是开的
（《空间的态射》引理 \ref{spaces-morphisms-lemma-fppf-open}），并且
$|V| = \bigcup_{i \in I} \bigcup_{j \in J_i} f_{ij}(|V_{ij}|)$
是 $|V|$ 的开覆盖。由于 $|V|$ 拟紧，该覆盖有有限加细。
设 $V_{i_1j_1}, \ldots, V_{i_Nj_N}$ 实现这一点。
则 $\{V_{i_kj_k} \to V\}_{k = 1, \ldots, N}$ 是 $V$ 的标准
fpqc 覆盖，并加细族
% P09-ERR-0335
$\{U_i \times_U V \to V\}$。证明完毕。
\end{proof}

\begin{lemma}
\label{spaces-topologies-lemma-fpqc}
设 $S$ 为概形，$X$ 为 $S$ 上的代数空间。
\begin{enumerate}
\item 若 $X' \to X$ 是同构，则 $\{X' \to X\}$ 是 $X$ 的 fpqc 覆盖。
\item 若 $\{X_i \to X\}_{i\in I}$ 是 fpqc 覆盖，并且对每个
$i$，$\{X_{ij} \to X_i\}_{j\in J_i}$ 是 fpqc 覆盖，则
$\{X_{ij} \to X\}_{i \in I, j\in J_i}$ 是 fpqc 覆盖。
\item 若 $\{X_i \to X\}_{i\in I}$ 是 fpqc 覆盖，而
$X' \to X$ 是代数空间的态射，则
$\{X' \times_X X_i \to X'\}_{i\in I}$ 是 fpqc 覆盖。
\end{enumerate}
\end{lemma}

\begin{proof}
(1) 是显然的。考虑 (3) 中的 $g : X' \to X$ 及 fpqc 覆盖
$\{X_i \to X\}_{i\in I}$。由《空间的态射》引理
\ref{spaces-morphisms-lemma-base-change-flat}，态射
$X' \times_X X_i \to X'$ 平坦。
若 $h' : Z \to X'$ 是从仿射概形到 $X'$ 的态射，则置
$h = g \circ h' : Z \to X$。关于 $\{X_i \to X\}_{i\in I}$
的假设意味着，存在标准 fpqc 覆盖
$\{Z_j \to Z\}_{j = 1, \ldots, n}$ 以及盖住 $h$ 的态射
$Z_j \to X_{i(j)}$，其中某些 $i(j) \in I$。
由纤维积的泛性质，还得到 $h'$ 上的态射
$Z_j \to X' \times_X X_{i(j)}$。
因此 $\{X' \times_X X_i \to X'\}_{i\in I}$ 是 fpqc 覆盖，
这就证明了 (3)。

\medskip\noindent
设 $\{X_i \to X\}_{i\in I}$ 与
$\{X_{ij} \to X_i\}_{j\in J_i}$ 如 (2) 所述。
设 $h : Z \to X$ 为从仿射概形到 $X$ 的态射。
由假设，存在标准 fpqc 覆盖
$\{Z_j \to Z\}_{j = 1, \ldots, n}$ 及盖住 $h$ 的态射
$h_j : Z_j \to X_{i(j)}$，其中某些 $i(j) \in I$。
仍由假设，存在标准 fpqc 覆盖
$\{Z_{j, l} \to Z_j\}_{l = 1, \ldots, n(j)}$
以及盖住 $h_j$ 的态射 $Z_{j, l} \to X_{i(j)j(l)}$，
其中某些 $j(l) \in J_{i(j)}$。
由《拓扑》引理 \ref{topologies-lemma-fpqc-affine-axioms}，
族 $\{Z_{j, l} \to Z\}$ 是标准 fpqc 覆盖。
因此 $\{X_{ij} \to X\}_{i \in I, j\in J_i}$ 是 fpqc 覆盖。
\end{proof}

\begin{lemma}
\label{spaces-topologies-lemma-recognize-fpqc-covering}
设 $S$ 为概形，$X$ 为 $S$ 上的代数空间。
假设 $\{f_i : X_i \to X\}_{i \in I}$ 是以 $X$ 为靶的代数空间态射族。
设 $U \to X$ 为从概形到 $X$ 的满平展态射。
则 $\{f_i : X_i \to X\}_{i \in I}$ 是 $X$ 的 fpqc 覆盖，
当且仅当 $\{U \times_X X_i \to U\}_{i \in I}$ 是 $U$ 的 fpqc 覆盖。
\end{lemma}

\begin{proof}
若 $\{X_i \to X\}_{i \in I}$ 是 fpqc 覆盖，则由引理
\ref{spaces-topologies-lemma-fpqc}，$\{U \times_X X_i \to U\}_{i \in I}$ 也是。
反之，假设 $\{U \times_X X_i \to U\}_{i \in I}$ 是 fpqc 覆盖。
设 $h : Z \to X$ 为从仿射概形到 $X$ 的态射。
则 $U \times_X Z \to Z$ 是概形的满平展态射，特别地是开态射。
因此可找出 $U \times_X Z$ 的有限多个仿射开集
$W_1, \ldots, W_t$，其像覆盖 $Z$。
对每个 $j$，把 $\{U \times_X X_i \to U\}_{i \in I}$
是 fpqc 覆盖这一条件应用于态射 $W_j \to U$，得到加细
$\{W_j \times_X X_i \to W_j\}_{i \in I}$ 的标准 fpqc 覆盖
$\{W_{jl} \to W_j\}$。
于是 $\{W_{jl} \to Z\}$ 是 $Z$ 的标准 fpqc 覆盖
（见《拓扑》引理 \ref{topologies-lemma-fpqc-affine-axioms}），
并加细
% P09-ERR-0336
$\{Z \times_X X_i \to X\}$，即得结论。
\end{proof}

\begin{lemma}
\label{spaces-topologies-lemma-refine-fpqc-schemes}
设 $S$ 为概形，$X$ 为 $S$ 上的代数空间。假设
$\mathcal{U} = \{f_i : X_i \to X\}_{i \in I}$ 是 $X$ 的 fpqc 覆盖。
则存在 $\mathcal{U}$ 的加细
$\mathcal{V} = \{g_i : T_i \to X\}$，它是 fpqc 覆盖且每个
$T_i$ 都是概形。
\end{lemma}

\begin{proof}
略。提示：对每个 $i$，取概形 $T_i$ 及满平展态射
$T_i \to X_i$，再验证 $\{T_i \to X\}$ 是 fpqc 覆盖。
\end{proof}

\noindent
待续……
